\documentclass[10pt,a4paper]{amsart}
\usepackage{amsmath,amssymb,amsthm,graphicx}
\usepackage{verbatim}
\newenvironment{reference}{\comment}{\endcomment}
%\newenvironment{reference}{}{}
\newenvironment{slogan}{\comment}{\endcomment}
\newenvironment{history}{\comment}{\endcomment}

% For commutative diagrams we use Xy-pic
\usepackage[all]{xy}

% We use 2cell for 2-commutative diagrams.
\xyoption{2cell}
% Theorem environments.
%
\theoremstyle{plain}
\newtheorem{theorem}[subsection]{Theorem}
\newtheorem{proposition}[subsection]{Proposition}
\newtheorem{lemma}[subsection]{Lemma}

\theoremstyle{definition}
\newtheorem{definition}[subsection]{Definition}
\newtheorem{example}[subsection]{Example}
\newtheorem{exercise}[subsection]{Exercise}
\newtheorem{situation}[subsection]{Situation}

\theoremstyle{remark}
\newtheorem{remark}[subsection]{Remark}
\newtheorem{remarks}[subsection]{Remarks}

\numberwithin{equation}{subsection}

% Macros
%
\def\lim{\mathop{\mathrm{lim}}\nolimits}
\def\colim{\mathop{\mathrm{colim}}\nolimits}
\def\Spec{\mathop{\mathrm{Spec}}}
\def\Hom{\mathop{\mathrm{Hom}}\nolimits}
\def\Ext{\mathop{\mathrm{Ext}}\nolimits}
\def\SheafHom{\mathop{\mathcal{H}\!\mathit{om}}\nolimits}
\def\SheafExt{\mathop{\mathcal{E}\!\mathit{xt}}\nolimits}
\def\Sch{\mathit{Sch}}
\def\Mor{\mathop{\mathrm{Mor}}\nolimits}
\def\Ob{\mathop{\mathrm{Ob}}\nolimits}
\def\Sh{\mathop{\mathit{Sh}}\nolimits}
\def\NL{\mathop{N\!L}\nolimits}
\def\CH{\mathop{\mathrm{CH}}\nolimits}
\def\proetale{{pro\text{-}\acute{e}tale}}
\def\etale{{\acute{e}tale}}
\def\QCoh{\mathit{QCoh}}
\def\Ker{\mathop{\mathrm{Ker}}}
\def\Im{\mathop{\mathrm{Im}}}
\def\Coker{\mathop{\mathrm{Coker}}}
\def\Coim{\mathop{\mathrm{Coim}}}

% Boxtimes
%
\DeclareMathSymbol{\boxtimes}{\mathbin}{AMSa}{"02}

%
% Macros for moduli stacks/spaces
%
\def\QCohstack{\mathcal{QC}\!\mathit{oh}}
\def\Cohstack{\mathcal{C}\!\mathit{oh}}
\def\Spacesstack{\mathcal{S}\!\mathit{paces}}
\def\Quotfunctor{\mathrm{Quot}}
\def\Hilbfunctor{\mathrm{Hilb}}
\def\Curvesstack{\mathcal{C}\!\mathit{urves}}
\def\Polarizedstack{\mathcal{P}\!\mathit{olarized}}
\def\Complexesstack{\mathcal{C}\!\mathit{omplexes}}
% \Pic is the operator that assigns to X its picard group, usage \Pic(X)
% \Picardstack_{X/B} denotes the Picard stack of X over B
% \Picardfunctor_{X/B} denotes the Picard functor of X over B
\def\Pic{\mathop{\mathrm{Pic}}\nolimits}
\def\Picardstack{\mathcal{P}\!\mathit{ic}}
\def\Picardfunctor{\mathrm{Pic}}
\def\Deformationcategory{\mathcal{D}\!\mathit{ef}}
\usepackage{cleveref}
\crefname{originallemma}{original lemma}{original lemmas}
\Crefname{originallemma}{Original lemma}{Original lemmas}
\makeatletter
\newlabel{lemma-pu-sheaf}{{00XK}{source}{}{sourcepointer-1}{}}
\newlabel{lemma-pu-sheaf@cref}{{[originallemma][0][]00XK}{[1][0][]source}{}{}{}}
\makeatother
\title{A cocontinuous continuous fully faithful functor gives a subtopos with a retraction}
\author{The Stacks Project Authors}
\date{stacks-project:sites.tex:local-snapshot-20261012; GNU FDL 1.2 or later}
\begin{document}
\maketitle
\noindent\textit{Editorial note.} All delivered mathematical text is the author's own, copied from the pinned Stacks Project chapter file sites-cohomology.tex; no statement or proof step was rewritten, reordered or removed.\par\medskip
\noindent\textit{Editorial note.} The proof is the author's complete argument; every step the author writes out is inside the counted span and no closing sentence is a placeholder.\par\medskip
\noindent\textit{Editorial note.} Every labelled reference occurring inside the counted spans is carried in assembly.external\_references as a non-load-bearing navigation pointer (same-chapter ones at their definition site, cross-chapter ones at the invocation inside the counted span), so no load-bearing mathematics is replaced by a stub. No bibliography citation occurs inside the counted spans.\par\medskip
\noindent\textit{Editorial note.} The checked chapter list for same-name collisions is the set of stacks chapters staged in this repository: chow, cohomology, divisors, etale-cohomology, flat, more-algebra, more-morphisms, morphisms, spaces-more-morphisms, varieties together with this card's chapter sites-cohomology.\par\medskip
\noindent\textit{Editorial note.} The preamble is the chapter's own retained minimal preamble (sources/stacks-sites-cohomology/retained-preamble.tex), a verbatim excerpt of the native preamble: the Xy-pic setup and the native theorem declarations with the mathematical macros.\par\medskip
\noindent\textit{Allowed standard prerequisites (editorial).}\begin{itemize}
\item The morphism of topoi associated to a cocontinuous and continuous functor. Same chapter, Lemma lemma-when-shriek, recorded as the context of this card. Supplies the functors g\_!, g\textasciicircum{}\{-1\}, g\_* and the description of the pullback.
\item Sheafification is not needed for the pushforward and the left adjoint. Same chapter, Lemmas lemma-pu-sheaf and lemma-when-shriek. Allows the composites to be computed on the level of presheaves.
\end{itemize}\medskip
\ifdefined\hypertarget\hypertarget{sourcepointer-1}{}\fi
\noindent\textit{Source reference.} Original lemma 00XK: sources/stacks-sites/pinned-original-source.tex, source lines 4897--4897. Navigation pointer invoked at proof line 4897; the counted span contains the pointer only and the cited result is not reproduced inside it, so the pointer is registered as a non-load-bearing navigation stub in the established form for Stacks material. No bibliography citation is involved.\par\medskip
\begin{lemma}
\label{lemma-when-shriek}
\begin{slogan}
Sheafification is redundant in topoi morphisms associated to simultaneously
continuous and cocontinuous site functors.
\end{slogan}
Let $\mathcal{C}$ and $\mathcal{D}$ be sites.
Let $u : \mathcal{C} \to \mathcal{D}$ be a functor.
Assume that
\begin{enumerate}
\item[(a)] $u$ is cocontinuous, and
\item[(b)] $u$ is continuous.
\end{enumerate}
Let $g : \Sh(\mathcal{C}) \to \Sh(\mathcal{D})$
be the associated morphism of topoi. Then
\begin{enumerate}
\item sheafification in the formula $g^{-1} = (u^p\ )^\#$ is
unnecessary, in other words $g^{-1}(\mathcal{G})(U) = \mathcal{G}(u(U))$,
\item $g^{-1}$ has a left adjoint $g_{!} = (u_p\ )^\#$, and
\item $g^{-1}$ commutes with arbitrary limits and colimits.
\end{enumerate}
\end{lemma}

\begin{lemma}
\label{lemma-back-and-forth}
Let $\mathcal{C}$ and $\mathcal{D}$ be sites.
Let $u : \mathcal{C} \to \mathcal{D}$ be a functor.
Assume that
\begin{enumerate}
\item[(a)] $u$ is cocontinuous,
\item[(b)] $u$ is continuous, and
\item[(c)] $u$ is fully faithful.
\end{enumerate}
For $g_!, g^{-1}, g_*$ as above
the canonical maps $\mathcal{F} \to g^{-1}g_!\mathcal{F}$
and $g^{-1}g_*\mathcal{F} \to \mathcal{F}$ are isomorphisms
for all sheaves $\mathcal{F}$ on $\mathcal{C}$.
\end{lemma}

\begin{proof}
Let $X$ be an object of $\mathcal{C}$.
In Lemmas \ref{lemma-pu-sheaf} and \ref{lemma-when-shriek} we have seen that
sheafification is not necessary for the functors
$g^{-1} = (u^p\ )^\#$ and $g_{*} = ({}_pu\ )^\#$.
We may compute
$(g^{-1}g_{*}\mathcal{F})(X) = g_{*}\mathcal{F}(u(X))
= \lim \mathcal{F}(Y)$. Here the limit
is over the category of pairs $(Y, u(Y) \to u(X))$
where the morphisms $u(Y) \to u(X)$ are not required to be
of the form $u(\alpha)$ with $\alpha$ a morphism of $\mathcal{C}$.
By assumption (c) we see that they automatically come from
morphisms of $\mathcal{C}$ and we deduce that the limit is the
value on $(X, u(\text{id}_X))$, i.e., $\mathcal{F}(X)$.
This proves that $g^{-1}g_{*}\mathcal{F} = \mathcal{F}$.

\medskip\noindent
On the other hand, $(g^{-1}g_{!}\mathcal{F})(X) =
g_{!}\mathcal{F}(u(X)) = (u_p\mathcal{F})^\#(u(X))$, and
$u_p\mathcal{F}(u(X)) = \colim \mathcal{F}(Y)$.
Here the colimit is over the category of pairs $(Y, u(X) \to u(Y))$
where the morphisms $u(X) \to u(Y)$ are not required to be
of the form $u(\alpha)$ with $\alpha$ a morphism of $\mathcal{C}$.
By assumption (c) we see that they automatically come
from morphisms of $\mathcal{C}$ and we deduce that the colimit is the
value on $(X, u(\text{id}_X))$, i.e., $\mathcal{F}(X)$. Thus for every
$X \in \Ob(\mathcal{C})$ we have
$u_p\mathcal{F}(u(X)) = \mathcal{F}(X)$.
Since $u$ is cocontinuous and continuous any covering of $u(X)$ in
$\mathcal{D}$ can be refined by a covering (!) $\{u(X_i) \to u(X)\}$
of $\mathcal{D}$ where $\{X_i \to X\}$ is a covering in $\mathcal{C}$.
This implies that $(u_p\mathcal{F})^+(u(X)) = \mathcal{F}(X)$ also,
since in the colimit defining the value of $(u_p\mathcal{F})^+$
on $u(X)$ we may restrict to the cofinal system of coverings
$\{u(X_i) \to u(X)\}$ as above. Hence we see that
$(u_p\mathcal{F})^+(u(X)) = \mathcal{F}(X)$ for all objects $X$
of $\mathcal{C}$ as well. Repeating this argument one more time
gives the equality $(u_p\mathcal{F})^\#(u(X)) = \mathcal{F}(X)$
for all objects $X$ of $\mathcal{C}$. This produces the desired
equality $g^{-1}g_!\mathcal{F} = \mathcal{F}$.
\end{proof}

\end{document}
